%-------------------------------%
%  Author: Alessandro Sciarra   %
%    Date: 11 Jul 2025          %
%-------------------------------%

% Partially switch to the green variant of the Z02 theme for this section
% NOTE: This is a low-level patch taking the colours from the
%       theme implementation. I need it only once in this session
%       and therefore I did not look for a more elegant way.
% NOTE: I did not switch completely to not interfere with the color-code
%       of sessions of the rest of the day!
\colorlet{PT}{UnrealFoodPills}
\colorlet{ALERT}{PT}
\colorlet{BGDARK}{GreenBlue}
\colorlet{BGLIGHT}{MicroscopeSlide}
\colorlet{TEXT}{Abyss}

%===================================================================================================%
\newcommand{\TextWithLink}[2]{\hyperlink{#2}{#1}}
\MakeSubsectionFramefalse
\newcommand{\MakeLabelledSubsectionWithCoverFrame}[2]%
{%
    \subsection{#2}
    \begin{frame}[plain, noframenumbering, label=#1]{}
        \subsectionpage
        % Add link to overview everywhere on the slide
        \begin{tikzpicture}[remember picture,overlay]
            %xelatex needs \XeTeXLinkBox, won't create a link unless it
            %finds text --- rules don't work without \XeTeXLinkBox.
            %Still builds correctly with pdflatex and lualatex
            \node at (current page.center){\hyperlink{TopicsWithLinks}{\XeTeXLinkBox{
                % See https://tex.stackexchange.com/a/60331 for nested tikzpicture
                \tikz{\node[minimum width=\paperwidth, minimum height=\paperheight]{};}
            }}};
        \end{tikzpicture}%
    \end{frame}
}
%Add a link to table of content on frames with a footline
\newif\ifAddLinkToGitHubWheel
\AddLinkToGitHubWheelfalse
\addtobeamertemplate{footline}{}{%
    \begin{tikzpicture}[remember picture,overlay]
        \ifAddLinkToGitHubWheel
            %xelatex needs \XeTeXLinkBox, won't create a link unless it
            %finds text --- rules don't work without \XeTeXLinkBox.
            %Still builds correctly with pdflatex and lualatex
            \node[anchor=north east, inner sep=2pt] at (current page.north east)
                {\hyperlink{TopicsWithLinks}{\XeTeXLinkBox{
\includegraphics[width=5mm]{Logo-GitHub}}}};
        \fi
    \end{tikzpicture}%
}%
%===================================================================================================%
\newlength\RootSize
\setlength\RootSize{30mm}
\begin{frame}[label=topics]{About GitHub}
    \hypertarget{TopicsWithLinks}{}
    % Define experience colors
    \colorlet{ACTION}{PB}
    \colorlet{BRANCH}{PT!50!yellow}
    \colorlet{ISSUE}{PT}
    \colorlet{FILE}{violet!75}
    \colorlet{PR}{PP}
    \colorlet{PROJECT}{PS}
    \colorlet{RELEASE}{PS!50!yellow}
    \colorlet{WIKI}{violet}
    \begin{center}
        \begin{tikzpicture}[remember picture, overlay]
            \begin{scope}[
                small mindmap,
                every node/.style={
                    concept, execute at begin node=\hskip0pt%, scale=0.9
                },
                root concept/.append style={
                    concept color=black, fill=white, line width=1.5mm, text width=\the\RootSize
                },
                text=white,
                actions/.style={concept color=ACTION,        faded/.style={concept color=ACTION!50}},
                branches/.style={concept color=BRANCH,       faded/.style={concept color=BRANCH!50}},
                files/.style={concept color=FILE,            faded/.style={concept color=FILE!50}},
                issues/.style={concept color=ISSUE,          faded/.style={concept color=ISSUE!50}},
                prs/.style={concept color=PR,                faded/.style={concept color=PR!50}},
                projects/.style={concept color=PROJECT,      faded/.style={concept color=PROJECT!50}},
                releases/.style={concept color=RELEASE,      faded/.style={concept color=RELEASE!50}},
                wiki/.style={concept color=WIKI,             faded/.style={concept color=WIKI!50}},
                grow cyclic,
                level 1/.append style={
                    level distance=35mm, font=\footnotesize, clockwise from=135, sibling angle=45, minimum size=18mm, text width=15mm
                }
                ]

                \node[root concept] (GH) at ($(current page.center)!0.18!(current page.east)$) {}
                child [branches] {
                    node {\TextWithLink{Branches}{BranchFrame}}
                }
                child [issues] {
                    node {\TextWithLink{Issues}{IssueFrame}}
                }
                child [prs] {
                    node {\TextWithLink{Pull\\ requests}{PrFrame}}
                }
                child [actions] {
                    node {\TextWithLink{Actions and CI}{ActionFrame}}
                }
                child [releases] {
                    node {\TextWithLink{Releases}{ReleaseFrame}}
                }
                child [projects] {
                    node {\TextWithLink{Projects}{ProjectFrame}}
                }
                child [wiki] {
                    node {\TextWithLink{Wiki and Webpage}{WikiFrame}}
                }
                child [files] {
                    node (SF) {\TextWithLink{Special\\ files}{FileFrame}}
                };
                \node [annotation, left=3mm, text width=15mm, fill=FILE!20, draw=none, text=FILE] at (SF.west)
                {
                    \begin{tabular}{@{}l}
                        $\circ\;$ README.md \\
                        $\circ\;$ LICENSE \\
                        $\circ\;$ CHANGELOG.md \\
                        $\circ\;$ AUTHORS.md \\
                        $\circ\;$ CONTRIBUTING.md \\
                        $\circ\;$ CITATION.cff \\
                    \end{tabular}
                };
            \end{scope}
            \node at (GH) {
\includegraphics[width=0.92\RootSize]{Logo-GitHub}};
        \end{tikzpicture}
    \end{center}
\end{frame}
%===================================================================================================%
\AddLinkToGitHubWheeltrue
\MakeLabelledSubsectionWithCoverFrame{BranchFrame}{Branches}
%===================================================================================================%
\begin{frame}{Branches}
    \begin{tikzpicture}[remember picture, overlay]
        \begin{scope}[scope on = <1>]
            \node (GF) at ($(current page.center)!0.2!(current page.east)$)
            {\includegraphics[height=\textheight]{Git-model}};
            \node[left = 5mm of GF, font=\large, text=PS] {
                \begin{tabular}{c}
                    Git-flow\\
                    branching\\
                    model\\
                \end{tabular}
            };
        \end{scope}
        \begin{scope}[scope on = <2>]
            \node (GT) at ([yshift=-7mm]current page.center)
            {\includegraphics[width=0.9\textwidth]{Git-trunk}};
            \node[above = 1mm of GT, font=\large, text=PS] {Trunk-based development};
        \end{scope}
    \end{tikzpicture}
    \begin{onlyenv}<3>
        \vspace{-6mm}
        \begin{varblock}{alert}[0.8\textwidth]{When does trunk-based development work best?}
            When you need to iterate quickly\\
            When you work mostly with senior developers
        \end{varblock}
        \begin{varblock}{example}[0.7\textwidth]{When does Git-flow work best?}
            When you have a lot of junior developers\\
            When you have an established product\\
            When you run an open-source project
        \end{varblock}
        \begin{varblock}{}[0.72\textwidth]{In academia}
            Development is rather slow and in small teams
        \end{varblock}
    \end{onlyenv}
    \PrepareURLsymbol[PB]
    \FrameRemark<3>{If you are interested in a nice comparison of the two development techniques, refer to \URL*{https://www.toptal.com/software/trunk-based-development-git-flow}{this nice overview}.}
\end{frame}
%===================================================================================================%
\MakeLabelledSubsectionWithCoverFrame{IssueFrame}{Issues}
%===================================================================================================%
\begin{frame}{Issues}{\alt<1>{A way to organise your (team's) work}{This is how it might look like}}
    \vspace{-6mm}
    \begin{varblock}{example}[0.85\textwidth]{When opening the issue}<only@1>
        \begin{enumerate}
            \item Choose a descriptive title
            \item Spend time on description!
            \item Assign somebody whenever possible
            \item Consistently use labels for issues properties
            \item Consistently use types to describe the issue entity
            \item Assign the issue to a project if reasonable
            \item Add the issue to a milestone if it makes sense
        \end{enumerate}
    \end{varblock}
    \begin{onlyenv}<2>
        \begin{center}
            \includegraphics[width=\textwidth]{GitHub-issues}
        \end{center}
    \end{onlyenv}
    \FrameRemark<1>{Consider issues relationships and associating a dedicate branch, if appropriate.}
\end{frame}
%===================================================================================================%
\MakeLabelledSubsectionWithCoverFrame{PrFrame}{Pull requests}
%===================================================================================================%
\begin{frame}{Pull requests}{A way to contribute to a repository}
    \vspace{-2mm}
    \begin{varblock}{example}[0.95\textwidth]{A clever merge procedure}
        \small
        A pull request is a way to prepare a merge of a Git branch into another one.
        Through a graphic interface, the changes can be reviewed and actions can be done.
    \end{varblock}
    \vspace{2mm}
    \begin{overlayarea}{\textwidth}{0.5\textheight}
        \begin{onlyenv}<1>
            \begin{center}
                \includegraphics[width=0.8\textwidth]{GitHub-PR}
            \end{center}
        \end{onlyenv}
        \begin{onlyenv}<2>
            \begin{itemize}
                \item Always assign at least one reviewer
                \item A PR can be opened as a \textbf{Draft}
                \item Use automatic linking to issues in the description
                \item \alert{Actions automatically start}
                \item The review is done in the ``Files changed'' tab
                \item \PB{The reviewer comments on code and can request changes}
            \end{itemize}
        \end{onlyenv}
    \end{overlayarea}
    \FrameRemark{Assignees are people in charge of merging the PR and may be who opened the PR, but also somebody else.}
\end{frame}
%~~~~~~~~~~~~~~~~~~~~~~~~~~~~~~~~~~~~~~~~~~~~%
\tikzset{
    space/.style={%
        thick, draw=#1, fill=#1!10, text=#1, rounded corners=1mm, font=\small, text width=14mm, align=center, minimum height=12mm
    }
}
\begin{frame}{PR reviewing process}{An example of typical workflow}
    \begin{tikzpicture}[overlay, remember picture]
        \coordinate (N0) at ($(current page.center)!1.0!(current page.north east)-(0.5,1)$);
        \foreach \n/\c [count=\i from 0, count=\ip from 1] in {Open PR/PP,
            Reviewer requests changes/PT,
            Discuss/PB,
            Implement changes/PQ}{
            \begin{scope}[scope on=<\ip->]
                \node[space=\c, below = 0mm of N\i, xshift=-25mm] (N\ip) {\n};
                \ifnum\i>0%
                \path[thick, to] (N\i.west) edge[out=180, in=90] (N\ip.north);
                \fi
            \end{scope}
        }
        \begin{scope}[scope on=<4->]
            \path[thick, draw, dotted] (N4.10)  -- ++(1.01,0) edge[out=0, in=270, looseness=1.0, to, dash phase=15pt] (N2.south);
        \end{scope}
        \begin{scope}[scope on=<5->]
            \node[thick, draw=PT, fill=PT!10, text=PT, diamond, font=\small,
            text width=14mm, align=center, inner sep=0mm, anchor=south] (NA) at (N1 |- N4.south) {Do actions pass?};
            \draw[thick, to] (N1.south) -- (NA.north);
            \draw[thick, to] (NA.west) -- ++(-0.3,0) node[pos=0.5, font=\footnotesize, anchor=south] {No} |- (N4.east);
            \path[thick, draw, dotted] (N4.350) -- ++(1.01,0) edge[out=0, in=270, looseness=0.6, to, dash phase=15pt] (NA.south);
        \end{scope}
        \begin{scope}[scope on=<6->]
            \node[space=PS, below = 0mm of N4, xshift=25mm] (N5) {Reviewer approves};
            \path[thick, to] (N4.south) edge[out=270, in=180] (N5.west);
            \draw[thick] (NA.east) -- ++(0.3,0) node[pos=0.5, font=\footnotesize, anchor=south] {Yes} coordinate (yes);
        \end{scope}
        \begin{scope}[scope on=<7->]
            \node[space=PP, below = 0mm of N5, xshift=33mm] (N6) {Merge PR};
            \path[thick, to] (N5.south) edge[out=270, in=180, looseness=0.8] (N6.west);
        \end{scope}
        \begin{scope}[scope on=<7->]
            \draw[thick, to] (yes) |- (N6.east);
        \end{scope}
    \end{tikzpicture}
\end{frame}
%===================================================================================================%
\MakeLabelledSubsectionWithCoverFrame{ActionFrame}{Actions and CI}
%===================================================================================================%
\begin{frame}[fragile]{Actions and continuous integration}{\only<1-2>{A world plenty of opportunities!}\only<3-4>{Some typical examples}\only<5>{Standard GitHub-hosted runners for \textbf{public} repositories}}
    \begin{onlyenv}<1>
        \begin{center}
            
\includegraphics[width=0.95\textwidth]{Meme-actions}
        \end{center}
    \end{onlyenv}
    \begin{onlyenv}<2>
        \begin{lstlisting}[style=MyBash]
            # List of all possible events that can trigger an action
            branch_protection_rule         page_build
            check_run                      public
            check_suite                    pull_request
            create                         pull_request_comment
            delete                         pull_request_review
            deployment                     pull_request_target
            deployment_status              push
            discussion                     registry_package
            discussion_comment             release
            fork                           repository_dispatch
            gollum                         schedule
            issue_comment                  status
            issues                         watch
            label                          workflow_call
            merge_group                    workflow_dispatch
            milestone                      workflow_run
        \end{lstlisting}
        \begin{center}
            Find out more about \URL[PB]{https://docs.github.com/en/actions/writing-workflows/choosing-when-your-workflow-runs/events-that-trigger-workflows}{events that trigger a workflow}!
        \end{center}
    \end{onlyenv}
    \begin{onlyenv}<3-4>
        \definecolor{myCyan}{cmyk}{1,0,0,0}
        \begin{varblock}{alert}[0.85\textwidth]{}<only@3>
            \alert{Actions use YAML files to configure what needs to be done!}
        \end{varblock}
        \begin{enumerate}
            \item<only@3> Run tests for default PR activity\\[-1mm]
            \Remark{The workflow runs when a pull request is opened or reopened or when the head branch of the pull request is updated}
            \begin{lstlisting}[style=MyBash, xleftmargin=1mm, xrightmargin=25mm, aboveskip=\medskipamount,
                               moredelim={[is][\color{myCyan}]{|-}{-|}},
                               moredelim={[is][\color{ForestGreen}]{|=}{=|}}]
                |-on-|:
                  |-pull_request-|:
                    |-branches-|:
                      - |=main=|
                      - |=release*=|
            \end{lstlisting}
            \item<only@3> Run tests when merged code into some branches
            \begin{lstlisting}[style=MyBash, xleftmargin=1mm, xrightmargin=25mm,
                               belowskip=-3mm, aboveskip=\medskipamount,
                               moredelim={[is][\color{myCyan}]{|-}{-|}},
                               moredelim={[is][\color{ForestGreen}]{|=}{=|}}]
                |-on-|:
                  |-push-|:
                    |-branches-|: [|=main=|, |=develop=|, |=release*=|]
            \end{lstlisting}
            \item<only@4> Selecting GitHub virtual machines and run something
            \begin{lstlisting}[style=MyBash, xleftmargin=1mm, xrightmargin=25mm,
                               belowskip=-6mm, aboveskip=1.5\medskipamount,
                               moredelim={[is][\color{myCyan}]{|-}{-|}},
                               moredelim={[is][\color{ForestGreen}]{|=}{=|}}]
                |-jobs-|:
                  |-Run-on-Ubuntu-|:
                    |-name-|: "Test on latest Ubuntu"
                    |-runs-on-|: |=ubuntu-latest=|
                    |-steps-|:
                      - |-uses-|: |=actions/checkout@v4=|
                      - |-name-|: "Print a greeting"
                        |-env-|:
                          |-MY_VAR-|: "Hi there"
                        |-run-|: |
                          |=echo "${MY_VAR}!"=|
            \end{lstlisting}
            \item<only@4> Using \URL[PB]{https://docs.github.com/en/actions/hosting-your-own-runners/managing-self-hosted-runners/about-self-hosted-runners}{a self-hosted runner}
            \begin{lstlisting}[style=MyBash, xleftmargin=1mm, xrightmargin=25mm,
                               belowskip=-8mm, aboveskip=1.5\medskipamount,
                               moredelim={[is][\color{myCyan}]{|-}{-|}},
                               moredelim={[is][\color{ForestGreen}]{|=}{=|}}]
                |-runs-on-|: [|=self-hosted=|, |=linux=|, |=x64=|]
            \end{lstlisting}
            \Remark{If you specify an array of strings or variables, your workflow will execute on any runner that matches all of the specified \,\texttt{runs-on}\, values.}
        \end{enumerate}
    \end{onlyenv}
    \begin{onlyenv}<5>
        \begin{center}
            \vspace{2mm}
            \includegraphics[height=0.7\textheight]{GitHub-actions-servers-public}
        \end{center}
    \end{onlyenv}
    \PrepareURLsymbol[PB]
    \FrameRemark<4-5>{Checkout \URL*{https://docs.github.com/en/actions/writing-workflows/workflow-syntax-for-github-actions}{the official documentation} to discover how to structure your actions. Check out existing actions, too!}
\end{frame}
%~~~~~~~~~~~~~~~~~~~~~~~~~~~~~~~~~~~~~~~~~~~~%
\begin{frame}[t]{Creating your first workflow}{Taken from \URL[PB]{https://docs.github.com/en/actions/writing-workflows/quickstart}{the official actions quick-start}}
    \vspace{1mm}
    \begin{overlayarea}{\textwidth}{0.75\textheight}
        \begin{onlyenv}<1>
            \quad\MakeEnumerateBox{1}~~Action files need to be placed in the \textbf{.github/workflows} folder\vphantom{\URL{}{}}
            \begin{center}
                \includegraphics[width=0.9\textwidth]{GitHub-actions-example-folder}
            \end{center}
        \end{onlyenv}
        \begin{onlyenv}<2>
            \quad\MakeEnumerateBox{2}~~Action files are written using \URL[PP]{https://docs.github.com/en/actions/writing-workflows/workflow-syntax-for-github-actions}{YAML workflows syntax}
            \begin{center}
                \includegraphics[height=0.6\textheight]{GitHub-actions-example-yaml}
            \end{center}
        \end{onlyenv}
        \begin{onlyenv}<3>
            \quad\MakeEnumerateBox{3}~~Actions are triggered by the event you chose in the YAML file\vphantom{\URL{}{}}
            \begin{center}
                \includegraphics[height=0.6\textheight]{GitHub-actions-example-UI}
            \end{center}
        \end{onlyenv}
        \begin{onlyenv}<4>
            \quad\MakeEnumerateBox{4}~~Actions log file is available on GitHub\vphantom{\URL{}{}}
            \begin{center}
                \includegraphics[height=0.6\textheight]{GitHub-actions-example-log}
            \end{center}
        \end{onlyenv}
    \end{overlayarea}
\end{frame}
%===================================================================================================%
\MakeLabelledSubsectionWithCoverFrame{ReleaseFrame}{Releases}
%===================================================================================================%
\begin{frame}{Releases}{An interface to associate metadata to a Git tag}
    \begin{center}
        \includegraphics[width=0.95\textwidth]{GitHub-release}
    \end{center}
\end{frame}
%===================================================================================================%
\MakeLabelledSubsectionWithCoverFrame{ProjectFrame}{Projects}
%===================================================================================================%
\begin{frame}{Projects}{Functionality at the Organisation or User level}
    \begin{center}
        \only<1>{\includegraphics[width=0.7\textwidth]{GitHub-projects-1}}
        \only<2>{\includegraphics[width=0.7\textwidth]{GitHub-projects-2}}
    \end{center}
\end{frame}
%===================================================================================================%
\MakeLabelledSubsectionWithCoverFrame{WikiFrame}{Wiki and webpage}
%===================================================================================================%
\begin{frame}{Wiki and webpage}
    \begin{onlyenv}<1>
        \begin{varblock}{example}[0.75\textwidth]{Wiki pages}
            You can enable and create Wiki pages associated to every repository.
            This is a separate repository with a dedicated UI bound to your main codebase.
        \end{varblock}
        \begin{center}
            \includegraphics[width=0.75\textwidth]{GitHub-wiki}
        \end{center}
    \end{onlyenv}
    \begin{onlyenv}<2>
        \PrepareURLsymbol[ALERT]
        \begin{varblock}{alert}[0.75\textwidth]{\URL*{https://pages.github.com}{GitHub pages}}
            You can create a website associated to your organisation/user or to any repository!
        \end{varblock}
        \begin{center}
            
\includegraphics[width=0.75\textwidth]{Fly}
        \end{center}
    \end{onlyenv}
    \PrepareURLsymbol[PB]
    \FrameRemark<2>{Checkout the \URL*{https://www.mkdocs.org}{MkDocs project} for possible inspiration.}
\end{frame}
%===================================================================================================%
\MakeLabelledSubsectionWithCoverFrame{FileFrame}{Special files}
%===================================================================================================%
\begin{frame}{Special files}
    \begin{description}[\hspace{35mm}]
        \item[\textbf{README.md}] An overview of your codebase
        \item[\textbf{LICENSE}] How is your material shipped to the public
        \item[\textbf{AUTHORS.md}] A list of people that contributed
        \item[\textbf{CONTRIBUTING.md}] How to become new developers
        \item[\textbf{CHANGELOG.md}] User-relevant changes between releases
        \item[\textbf{CITATION.cff}] How to cite your software
    \end{description}
    \vspace{3mm}
    \PrepareURLsymbol[PB]
    \begin{varblock}{alert}[0.68\textwidth]{Take advantage of this!}
        Most files usually use \URL*{https://docs.github.com/en/get-started/writing-on-github/getting-started-with-writing-and-formatting-on-github/basic-writing-and-formatting-syntax}{Markdown formatting}
    \end{varblock}
\end{frame}
%===================================================================================================%
\MakeSubsectionFrametrue
\AddLinkToGitHubWheelfalse
% Restore blue colours of Z02 theme for other sections
\colorlet{PB}{Marine}
\colorlet{PP}{SkyBlue}
\colorlet{PS}{BladeOfGrass}
\colorlet{PT}{CoralRed}
\colorlet{PQ}{Schizophrenia}
\colorlet{ALERT}{UnrealFoodPills}
\colorlet{BGDARK}{Aoi}
\colorlet{BGLIGHT}{BeachStorm!25}
\colorlet{TEXT}{Abyss}